% =====================================================================
% TEMPLATE — Problem description (methodology asset A-P1-tpl)
% =====================================================================
% Purpose: formalize a problem so any technique (heuristic, MILP, ...) can be
% built on top of it. Fill every <PLACEHOLDER>; delete guidance comments.
%
% Structure (do not change the skeleton):
%   Part 1 — Problem presentation (narrative: context, physical/logical setup,
%            how the cost/metric arises, key assumptions, operating regime).
%   Part 2 — Problem statement (explicit: Decision, Requirements, Objective,
%            Input information).
%
% Meant to be \input from a main.tex and compiled with a compile.py; see the
% working example harness in docs/<problem>/ (main.tex + compile.py).
% =====================================================================

% =====================================================================
\part{Problem presentation}
% =====================================================================

\section{Overview}
% 2–4 sentences: what system/process, what we optimize, and the single
% headline metric. State the central objective explicitly.
<OVERVIEW: what the system does, what is optimized, and the headline metric>

\section{Setup}
% The physical or logical building blocks and how items/entities flow through
% them. Use a numbered list of components, each with a short bullet list of its
% defining properties and assumptions.
\begin{enumerate}
  \item \textbf{<COMPONENT 1>.}
    \begin{itemize}
      \item <property / assumption>
    \end{itemize}
  \item \textbf{<COMPONENT 2>.} <...>
\end{enumerate}

\section{How the cost/metric arises}
% Explain, mechanically, what drives the objective metric up or down — the lever
% the solution will pull. This is the intuition the technique must exploit.
\begin{itemize}
  \item <when is a unit of cost incurred?>
  \item <what consolidates / avoids cost?>
\end{itemize}

\subsection{Key assumptions}
% Every simplifying assumption that changes the model. Be explicit: what is NOT
% tracked, what is assumed infinite/instant/feasible, what is out of scope.
\begin{itemize}
  \item <assumption 1>
\end{itemize}

% OPTIONAL — include only if the problem has distinct operating regimes /
% variants (e.g. a simplified baseline vs the target extension). Contrast them.
\section{Operating regime / variant}
\begin{itemize}
  \item \textbf{<Baseline / simplified>.} <description>
  \item \textbf{<Target / extension>.} <description and why it is harder/realistic>
\end{itemize}

% =====================================================================
\part{Problem statement}
% =====================================================================

\section{Decision}
% What the solver decides. One enumerated item per decision. If decisions are
% coupled, say so explicitly after the list.
\begin{enumerate}
  \item \textbf{<Decision 1>} --- <what is chosen>.
  \item \textbf{<Decision 2>} --- <what is chosen>.
\end{enumerate}

\section{Requirements}
% Hard constraints any feasible solution must satisfy. Group related constraints.
% Also state explicitly anything that is NOT a requirement (to bound the model).
%
% STABLE-ID CONVENTION (keep it): tag each requirement with a durable ID
% [RQ01], [RQ02], ... assigned SEQUENTIALLY as requirements are added. An ID is
% NEVER reused and NEVER renumbered: if a requirement is dropped, retire its ID
% (leave a gap) rather than reassigning it; new requirements always take the
% next unused number. These IDs are referenced downstream — the solution checker
% cites them to report which requirement a solution violates, and MILP
% constraints may cite them (e.g. a \label comment) for traceability.
\begin{itemize}
  \item \textbf{[RQ01]} <requirement>
  \item \textbf{[RQ02]} <requirement>
\end{itemize}

\section{Objective}
% One sentence: minimize/maximize <quantity> subject to the requirements.
\textbf{<Minimize/Maximize> <the metric>} subject to the requirements above.

\section{Input information}
% What defines an instance. One bullet per input entity with its fields.
\begin{itemize}
  \item \textbf{<Input 1>.} <fields / meaning>
  \item \textbf{<Input 2>.} <fields / meaning>
\end{itemize}

% OPTIONAL — how instances will be varied to produce scenarios / KPIs.
\medskip
\noindent\textbf{Scenario evaluation.} <which parameters are swept and why>
